\begin{table*}
\begin{tabular}{@{}cp{6cm}cccccc@{}}
\toprule
\textbf{Map Key} & \textbf{County, State (City Within County)} & \textbf{RFI$^\mathbf{\ast}$} & \textbf{Population} & \textbf{RPP$^\mathbf{\dag}$} & \textbf{RPA$^\mathbf{\ddag}$} & \textbf{Homelessness$^\mathbf{\nabla}$} & \textbf{GINI$^\mathbf{\times}$} \\ \midrule \addlinespace[3pt]
\multicolumn{8}{c}{\textbf{Counties / Cities Comparable to San Francisco County (San Francisco, CA, USA)}} \\ \addlinespace[4pt]
1 & San Francisco County, California (San Francisco) & 0.75 & 851,036 & 1032 & 131 & 98 & 0.52 \\
2 & Multnomah County, Oregon (Portland) & 0.56 & 808,098 & 1198 & 237 & 91 & 0.47 \\
3 & Erie County, New York (Buffalo) & 0.47 & 951,232 & 1342 & 134 & 59.50 & 0.463 \\
4 & Baltimore County, Maryland (Baltimore) & 0.63 & 850,737 & 997 & 99 & 6.64 & 0.456 \\
5 & El Paso County, Texas (El Paso) & 0.69 & 863,832 & 1919 & 99 & 10.57 & 0.466 \\ \midrule \addlinespace[4pt]
\multicolumn{8}{c}{\textbf{Counties / Cities Comparable to St. Joseph County (South Bend, IN, USA)}} \\ \addlinespace[4pt]
A & St. Joseph County, Indiana (South Bend) & 0.52 & 272,388 & 1378 & 97 & 8.00 & 0.47 \\
B & Winnebago County, Illinois (Rockford) & 0.57 & 284,591 & 1583 & 134 & 28.74 & 0.452 \\
C & Kalamazoo County, Michigan (Kalamazoo) & 0.43 & 261,426 & 1297 & 83 & 25.00 & 0.46 \\
D & Lackawanna County, Pennsylvania (Scranton) & 0.38 & 215,672 & 1252 & 238 & 8.40 & 0.456 \\
E & Washington County, Arkansas () & 0.60 & 247,331 & 1466 & 80 & 32.14 & 0.483 \\ \bottomrule
\end{tabular}
{\raggedright
$^\mathrm{\ast}$RFI: Racial Fractionalization Index \\
$^\mathrm{\dag}$RPP: Rate of People Below Poverty Line (per 10k) \\
$^\mathrm{\ddag}$RPA: Rate of People With Public Assistance (per 10k)\\
$^\mathrm{\nabla}$Homelessness: Homelessness Rate (per 10k) \\
$^\mathrm{\times}$GINI: Income Inequality (GINI)\\
}
\end{table*}